% ----------------------------------------------------------
\chapter{Discussões}\label{cap:discussoes}
% ----------------------------------------------------------

\section{Aplicabilidade na Atenção Primária}\label{sec:disc-aplicabilidade}
A regra clínica avaliada encaminhou 84,0\% do conjunto de teste a exame, em um ponto de operação único e não ajustável, de maneira que, aplicada à população atendida pela APS, a regra pouco seleciona quem examinar, deslocando para a etapa laboratorial praticamente toda a demanda. O modelo de referência, por sua vez, permite escolher o ponto de operação conforme a capacidade da rede, característica relevante em um contexto de atendimento predominantemente por demanda espontânea e com capacidade limitada \cite{geremia2020atenccao}. No ponto de operação principal, o modelo encaminhou 40,5\% do conjunto de teste, reduzindo os encaminhamentos em 736, isto é, em 51,8\%, ao custo de 19 casos não detectados em relação à regra. Além disso, a precisão, que corresponde à proporção de encaminhados que de fato apresentaram o desfecho, passou de 0,081 na regra para 0,140 no modelo.

Essa forma de avaliar uma estratégia de triagem pela fração da população que ela encaminha encontra paralelo na literatura nacional. Ao apresentar o BrDMrisc, \textcite{bracco2026brdmrisc} destacam como propriedade favorável do escore a identificação de uma fração mais manejável da população como de alto risco, de 20,0\% contra 40,6\% da estratégia baseada apenas na glicemia de jejum alterada, com probabilidade de desenvolver diabetes entre os identificados de 28,5\% contra 17,1\%. Cabe ressaltar que o BrDMrisc estima o risco de diabetes incidente a partir de uma coorte longitudinal, de maneira que os seus valores não são comparáveis aos deste trabalho, sendo a convergência restrita ao critério de avaliação da triagem.

No que se refere à capacidade discriminativa, a revisão sistemática de \textcite{CarrilloLarco2020} reuniu quatro estudos realizados na América do Sul que reestimaram o FINDRISC para a identificação de diabetes não diagnosticado, com prevalências entre 2,6\% e 5,1\%, e reportou valores de AUC-ROC entre 0,68 e 0,75. Os modelos deste trabalho, com AUC-ROC entre 0,770 e 0,786 no conjunto de teste, situam-se na mesma ordem de grandeza. Essa aproximação, no entanto, não constitui uma comparação direta, visto que os estudos revisados utilizaram amostras locais, um desfecho definido pela combinação de glicemia de jejum, HbA1c e teste de tolerância à glicose, restrito a casos não diagnosticados, além de variáveis comportamentais e de histórico familiar excluídas deste trabalho, enquanto o desfecho aqui adotado corresponde à HbA1c $\geq$ 6,5\% em uma subamostra nacional, sem distinção entre casos previamente diagnosticados e não diagnosticados. Os autores apontaram ainda que todos os estudos apresentaram alto risco de viés e que nenhum reportou métricas de calibração, aspecto avaliado neste trabalho e cujo resultado, apresentado no capítulo de Resultados, restringe a leitura das saídas dos modelos a escores de ordenação.

\section{O Critério de Sucesso no Regime de Sensibilidade Equiparada}\label{sec:disc-criterio}
O modelo de referência não cumpriu o critério de sucesso estabelecido antes da execução, e esse resultado foi mantido como tal, sem redefinição do critério, do regime de comparação ou do modelo de referência. A análise dos números que o compõem, entretanto, qualifica o seu alcance. No regime de sensibilidade equiparada, o modelo de referência encaminhou 40 indivíduos a mais que a regra para detectar 2 casos a mais, isto é, 20 encaminhamentos por caso adicional, valor superior ao NNS da própria regra, de 12,36, o que elevou o seu NNS acima do dela. O \textit{Random Forest} com o conjunto A, por sua vez, encaminhou 36 indivíduos a mais para detectar 3 casos a mais, isto é, 12 encaminhamentos por caso adicional, valor ligeiramente inferior ao NNS da regra, e cumpriu o critério por uma diferença de 0,009.

Essa diferença depende de um único caso, visto que, com 117 casos detectados em vez de 118, o NNS dessa combinação seria de 12,45, superior ao da regra. Diante disso, a substituição do modelo de referência por essa combinação foi descartada por dois motivos. O primeiro é metodológico, visto que a substituição constituiria a escolha do modelo a partir do resultado da própria avaliação, contrariando a regra definida antes da execução. O segundo decorre dos próprios números, visto que a substituição trocaria uma escolha fundamentada em regra prévia por uma vantagem da ordem de um caso, em uma partição de teste com 119 casos positivos.

Além disso, o critério foi decidido em uma região de operação na qual todas as estratégias encaminharam entre 84,0\% e 87,5\% do conjunto de teste e detectaram entre 115 e 118 dos 119 casos positivos, de maneira que as diferenças entre elas se resumiram a poucos casos e a algumas dezenas de encaminhamentos. Essa região decorre da própria regra clínica, que, ao encaminhar todo indivíduo com idade igual ou superior a 35 anos ou com excesso de peso, abrange a maior parte da amostra avaliada, de maneira que equiparar a sensibilidade da regra exige que os modelos operem no extremo de sua curva, onde a ordenação dos indivíduos oferece pouca margem de seleção. Por consequência, o regime equiparado responde se os modelos superam a regra no ponto em que ela opera, porém esse ponto corresponde a uma triagem que encaminha quase toda a população avaliada, situação distinta da que motivou o trabalho, cuja proposta reside na priorização dos encaminhamentos sob capacidade limitada.

Diante disso, o resultado negativo não demonstra que os modelos sejam inferiores à regra, tampouco que a superem, mas sim que, na região de operação da regra, as estratégias se mostraram praticamente equivalentes, com diferenças que dependem de um ou dois casos. A vantagem dos modelos, quando existe, situa-se na possibilidade de operar em outros pontos da curva, conforme discutido na seção de Aplicabilidade na Atenção Primária, argumento que, no entanto, não foi formalizado como critério de sucesso antes da execução e, por esse motivo, não substitui o critério estabelecido.

\section{\textit{Random Forest} e Regressão Logística}\label{sec:disc-algoritmos}
% [parágrafo sobre a PR-AUC: aguarda a decisão sobre a decomposição por faixa]

Parte da vantagem do \textit{Random Forest} decorre do comportamento da Regressão Logística no conjunto de teste. Na validação cruzada, os cinco modelos atingiram sensibilidade de 0,802, porém, no teste, o \textit{Random Forest} manteve valores de 0,782 e 0,807, enquanto a Regressão Logística apresentou 0,748 e 0,765. No conjunto B, essa queda de 0,037 corresponde a mais de quatro vezes o desvio padrão da própria regressão entre as partições, de 0,008, de maneira que o resultado do teste se afasta do comportamento observado na validação cruzada. Como a avaliação foi realizada em uma única partição de teste, com 119 casos positivos, na qual cada caso corresponde a 0,008 de sensibilidade, não é possível distinguir se essa queda reflete uma menor capacidade de generalização do ponto de operação da regressão ou uma particularidade da partição.

O teste pareado reforça essa leitura cautelosa, visto que a diferença de acerto global entre os modelos não foi significativa, enquanto a floresta detectou cinco casos a mais, incluindo todos os detectados pela regressão, e encaminhou 19 indivíduos sem o desfecho a mais. Em termos operacionais, esses cinco casos adicionais custaram 4,8 encaminhamentos cada, valor inferior ao NNS de ambos os modelos, e, no regime de sensibilidade equiparada, a floresta encaminhou menos indivíduos e detectou ao menos o mesmo número de casos nos dois conjuntos. Diante disso, a escolha do \textit{Random Forest} como modelo de referência é sustentada pelos indicadores operacionais, porém a magnitude da vantagem, de cinco casos em uma partição com 119 positivos, recomenda que ela seja confirmada em partições ou bases de dados adicionais. Além disso, a análise realizada não permite atribuir a diferença a um mecanismo específico, como a captura de interações entre variáveis, visto que nenhum experimento foi desenhado para isolar esse efeito.

\section{Desbalanceamento e Calibração}\label{sec:disc-calibracao}
% por que reponderar se o limiar já trata o desbalanceamento (van den Goorbergh)
% mecanismos distintos: logística (deslocamento compatível com o reponderamento) x floresta (média de votos e reponderamento por subamostra)
% escore de ordenação, não risco; recalibração posterior (Platt ou isotônica) como trabalho futuro

\section{Multicolinearidade e Explicabilidade}\label{sec:disc-explicabilidade}
% VIF não vê redundância por ponto de corte; escolha de A e B por interpretabilidade
% predição confirmada só no par de pressões; hipóteses: correlação de 0,79 não é redundância; subamostragem de colunas reparte o crédito
% topo estável entre regimes; restrição de leitura das contribuições negativas

\section{Limitações}\label{sec:limitacoes}
% métricas amostrais, não ponderadas; partição única; desenho transversal (associação, não incidência)
% calibração; SHAP da logística não executado; hipótese de ausência não aleatória do colesterol não testada
% IC bootstrap de raça/cor não realizado; regra clínica simplificada (sem histórico familiar)
% protótipo da fase de implantação, se não concluído
